\chapter{Theoretical Background}\label{cha:background}

\section{Neural Tangent Kernel}\label{sec:ntk}

As discussed in Chapter~\ref{cha:intro}, for nonlinear neural networks the
gradient flow dynamics are parameter-dependent and do not admit a simple
closed-form description. One approach to recovering a tractable analytical
framework is to consider regimes in which the network behaves approximately
linearly around its random initialization. The \emph{Neural Tangent Kernel}
(NTK), introduced by \citet{jacot2018ntk}, formalizes this idea by studying
training dynamics through a first-order linearization in parameter space,
leading to a kernel-based description of learning in wide neural networks.

\paragraph{Linearization around initialization}

Let $f(x;\theta)$ denote a neural network with parameters
$\theta \in \mathbb{R}^p$, initialized at $\theta_0$. A first-order Taylor
expansion around $\theta_0$ yields
\begin{equation}
	f(x;\theta)
	\;\approx\;
	f(x;\theta_0)
	+
	\nabla_\theta f(x;\theta_0)^\top (\theta - \theta_0),
	\label{eq:ntk-linearization}
\end{equation}
where $f(x;\theta_0)$ is the network output at initialization and
\[
	\phi(x) := \nabla_\theta f(x;\theta_0)
\]
defines the tangent feature map. Locally around initialization, the network
thus behaves as a linear model in parameter space,
\begin{equation}
	f(x;\theta) \approx f(x;\theta_0) + \phi(x)^\top (\theta - \theta_0).
	\label{eq:ntk-linear-model}
\end{equation}

\begin{definition}[Neural Tangent Kernel \citep{jacot2018ntk}]
	Given an initialization $\theta_0$, the neural tangent kernel is defined as
	\[
		\Theta_0(x,x')
		:= \nabla_\theta f(x;\theta_0)^\top
		\nabla_\theta f(x';\theta_0).
	\]
\end{definition}

The NTK measures how infinitesimal parameter updates couple the network outputs
at different inputs and plays the role of an empirical covariance operator in
the tangent feature space.

\paragraph{Training dynamics induced by the NTK}

We consider training with the squared loss using gradient descent with step size $\eta>0$,
\[
	L(\theta) = \frac{1}{2}\sum_{i=1}^n \bigl(f(x_i;\theta) - y_i\bigr)^2,
	\qquad \theta_{t+1} = \theta_t - \eta \nabla_\theta L(\theta_t),
	\qquad t = 0,1,2,\dots
\]
Define $f_t(x) := f(x;\theta_t)$. The gradient of the loss is
\[
	\nabla_\theta L(\theta_t)
	=
	\sum_{j=1}^n \bigl(f_t(x_j)-y_j\bigr)\,\nabla_\theta f(x_j;\theta_t).
\]
Applying a first-order Taylor expansion of $f(\cdot;\theta)$ around $\theta_t$
gives
\[
	f_{t+1}(x_i)
	\approx
	f_t(x_i)
	+
	\nabla_\theta f(x_i;\theta_t)^\top(\theta_{t+1}-\theta_t).
\]
Substituting the gradient descent update
$\theta_{t+1}-\theta_t = -\eta \nabla_\theta L(\theta_t)$
yields
\[
	f_{t+1}(x_i)
	=
	f_t(x_i)
	-
	\eta
	\sum_{j=1}^n
	\nabla_\theta f(x_i;\theta_t)^\top
	\nabla_\theta f(x_j;\theta_t)\,
	\bigl(f_t(x_j)-y_j\bigr),
\]
which can be written compactly as
\begin{equation}
	f_{t+1}(x_i)
	=
	f_t(x_i)
	-
	\eta
	\sum_{j=1}^n
	\Theta_t(x_i,x_j)\bigl(f_t(x_j)-y_j\bigr),
	\label{eq:ntk-prediction-dynamics}
\end{equation}
where the iteration-dependent NTK is
\[
	\Theta_t(x_i,x_j)
	:=
	\nabla_\theta f(x_i;\theta_t)^\top
	\nabla_\theta f(x_j;\theta_t).
\]
If we stack predictions into
$f_t := (f_t(x_1),\dots,f_t(x_n))^T \in \mathbb{R}^n$, in vector form, the dynamics read
\begin{equation}
	f_{t+1} = f_t - \eta\,\Theta_t (f_t - y).
	\label{eq:ntk-vector-dynamics}
\end{equation}
where $
	\bar\Theta
	=
	\bigl(\bar\Theta(x_i,x_j)\bigr)_{i,j=1}^n
	\in \mathbb{R}^{n\times n}$ is the kernel gram matrix.
In general, $\Theta_t$ evolves during training, reflecting changes in the
network’s tangent features.

\paragraph{One-hidden-layer NTK and infinite-width limit}

To make the NTK explicit, consider a two-layer network of width $m$,
\begin{equation}
	f(x)
	=
	\frac{1}{\sqrt{m}}
	\sum_{\alpha=1}^m
	v_\alpha\,\varphi(w_\alpha^\top x),
	\qquad
	v_\alpha \in \mathbb{R}, \quad w_\alpha, x \in \mathbb{R}^d.
	\label{eq:1hl-network}
\end{equation}
Each neuron is parameterized by $\theta_\alpha = (v_\alpha,w_\alpha)$ and
randomly initialized as
\[
	v_\alpha \sim \mathcal{N}(0,\sigma_v^2),
	\qquad
	w_\alpha \sim \mathcal{N}\!\left(0,\sigma_w^2I_d\right),
\]
independently across $\alpha = 1,\dots,m$. $\varphi: \mathbb{R} \to \mathbb{R}$ is a non-linear activation function.

\paragraph{Finite-width NTK at initialization.}
By direct computation,
\begin{align}
	\nabla_{v_\alpha} f(x)
	 & = \frac{1}{\sqrt{m}}\varphi(w_\alpha^\top x),                 \\
	\nabla_{w_\alpha} f(x)
	 & = \frac{1}{\sqrt{m}}\,v_\alpha\,\varphi'(w_\alpha^\top x)\,x.
\end{align}
Substituting into the definition of the NTK yields the empirical kernel
\begin{equation}
	\Theta_0^{(m)}(x,x')
	=
	\frac{1}{m}\sum_{\alpha=1}^m
	\varphi(w_\alpha^\top x)\,\varphi(w_\alpha^\top x')
	+
	\frac{1}{m}\sum_{\alpha=1}^m
	v_\alpha^2\,
	\varphi'(w_\alpha^\top x)\,
	\varphi'(w_\alpha^\top x')\,
	x^\top x'.
	\label{eq:ntk-1hl-finite}
\end{equation}

\paragraph{Infinite-width limit.}
Each sum in \eqref{eq:ntk-1hl-finite} is an empirical average of i.i.d.\ terms.
By the strong law of large numbers,
\[
	\Theta_0^{(m)}(x,x')
	\;\xrightarrow{\text{a.s.}}\;
	\bar\Theta(x,x')
	\qquad \text{as } m \to \infty,
\]
where the limiting NTK is deterministic and given by
\begin{equation}
	\bar\Theta(x,x')
	=
	\mathbb{E}_{w}\!\big[\varphi(w^\top x)\,\varphi(w^\top x')\big]
	+
	\sigma_v^2\,x^\top x'\,
	\mathbb{E}_{w}\!\big[\varphi'(w^\top x)\,\varphi'(w^\top x')\big].
	\label{eq:ntk-1hl-infinite}
\end{equation}
This argument extends to deep fully connected networks, where both the forward
covariance kernel and the NTK satisfy recursive layerwise equations in the
infinite-width limit \citep{jacot2018ntk, lee2020wide}.

\subsection{Constant-kernel regime and spectral dynamics}\label{sec:ntk-spectral-dynamics}

In the infinite-width limit, the NTK remains constant throughout training,
$\Theta_t \equiv \bar\Theta$. Equation~\eqref{eq:ntk-vector-dynamics} then reduces
to the linear iteration
\begin{equation}
	f_{t+1} = f_t - \eta\,\bar\Theta (f_t - y).
	\label{eq:ntk-constant-kernel}
\end{equation}
Letting $r_t := f_t - y$, we obtain
\[
	r_{t+1} = (I - \eta \bar\Theta)\,r_t.
\]
Assuming $\eta < 2/\lambda_{\max}(\bar\Theta)$, the residual admits the closed
form
\begin{equation}
	r_t = (I - \eta \bar\Theta)^t r_0,
	\qquad
	f_t = y + (I - \eta \bar\Theta)^t (f_0 - y).
	\label{eq:ntk-solution}
\end{equation}
As in the linear setting of Chapter~\ref{cha:intro}, convergence occurs
independently along the eigenvectors of $\bar\Theta$, with geometric rates
determined by the corresponding eigenvalues.

\paragraph{NTK dynamics, solution structure, and implicit bias}

In the constant-kernel (infinite-width) regime, the NTK prediction dynamics
reduce to a linear iteration in prediction space. Eq \ref{eq:ntk-constant-kernel},
admits fixed points $f^\star$ satisfying
\begin{equation}
	\bar\Theta\,(f^\star - y) = 0.
	\label{eq:ntk-fixed-point}
\end{equation}
Equation~\eqref{eq:ntk-fixed-point} characterizes the set of global minimizers of
the squared loss in the NTK regime. As in linear least squares, the structure of
this solution set depends on the rank of the operator $\bar\Theta$.

If $\bar\Theta$ is strictly positive definite on the training set, the solution
is unique and satisfies $f^\star = y$. When $\bar\Theta$ is singular, the loss
admits infinitely many interpolating solutions in prediction space, differing
by elements of $\ker(\bar\Theta)$. In this case, gradient descent converges to a
particular fixed point determined by the initialization.

\paragraph{Tangent features and operator factorization.}
To formulate the NTK in a way that remains meaningful in the infinite-width
limit, it is convenient to view parameter perturbations as elements of a
Hilbert space $\mathcal{H}_\theta$ equipped with an inner product
$\langle\cdot,\cdot\rangle$. In the finite-dimensional case,
$\mathcal{H}_\theta=\mathbb{R}^p$ with the Euclidean inner product; in the
infinite-width limit, $\mathcal{H}_\theta$ denotes the corresponding limit
space of parameter perturbations, often referred to as the tangent parameter
space \citep{jacot2018ntk, lee2020wide}.

In this setting, the limiting NTK Gram matrix on the training set can be
expressed as
\begin{equation}
	\bar\Theta = J_0 J_0^\ast,
\end{equation}
where $J_0 : \mathcal{H}_\theta \to \mathbb{R}^n$ denotes the Jacobian of the
network outputs with respect to parameters at initialization, viewed as a
linear operator,
\[
	(J_0 h)_i = \langle \nabla_\theta f(x_i;\theta_0),\,h\rangle,
\]
and $J_0^\ast$ is its adjoint. This representation makes explicit that
$\bar\Theta$ is symmetric and positive semidefinite, and that it acts as an
empirical covariance operator for the tangent features induced by the network
at initialization \citep{jacot2018ntk}.

\paragraph{Implicit bias and representer viewpoint.}
When $\bar\Theta$ is singular, the Moore--Penrose pseudoinverse $\bar\Theta^+$
acts as the inverse on $\mathrm{range}(\bar\Theta)$ and annihilates
$\ker(\bar\Theta)$. Consequently, constant-kernel NTK training eliminates only
the residual component lying in $\mathrm{range}(\bar\Theta)$ while preserving
the component in $\ker(\bar\Theta)$. This behavior reflects an implicit
regularization effect analogous to linear least squares and can be interpreted
as convergence to a minimum-norm solution in the reproducing kernel space
associated with $\bar\Theta$, as formalized by representer-theorem results
\citep{bietti2019inductivebiasneuraltangent,
	bartolucci2021understandingneuralnetworksreproducing}.

\paragraph{Limitations of stable kernels and feature learning.}
Although the NTK framework provides a clean and tractable description of
training dynamics, the stability of the kernel also introduces inherent
limitations. When the kernel does not change during training, learning is
effectively confined to a fixed feature space. Indeed, any positive
semidefinite kernel admits a representation
$K(x,x')=\langle \Phi(x),\Phi(x')\rangle_{\mathcal H}$, so that optimization in
the NTK regime corresponds to fitting a linear model in the associated
reproducing kernel Hilbert space $\mathcal H$.

In contrast, for finite-width networks the NTK typically evolves during
training, implying that the tangent features
$\Phi_t(x)=\nabla_\theta f(x;\theta_t)$ also change over time. This evolution
allows the network to adapt its representation to the data, a phenomenon
commonly referred to as \emph{feature learning}. Such effects are absent in a
strictly stable-kernel regime, where only the coefficients of fixed features
are adjusted.

The importance of feature learning is already visible in simplified settings.
Even purely linear networks exhibit nontrivial training dynamics due to the
coupling of parameters across layers, leading to implicit biases that cannot
be captured by a fixed kernel model
\citep{saxe2014exactsolutionsnonlineardynamics, tu2024mixeddynamicslinearnetworks}.
From this viewpoint, the NTK regime can be seen as a useful but restrictive
approximation, which motivates studying regimes where kernel evolution and
feature learning play an explicit role.


\section{Spectral dynamics on the circle}
\label{sec:background-spectral-dynamics-circle}

The constant-kernel NTK regime gives a linear description of training in
function space. We now specialize this viewpoint to the setting used throughout
the thesis: inputs on the unit circle \(S^1\). This gives a clean reference
model in which the operator governing training is diagonalized by Fourier modes.

Let \(S^1\) be parameterized by \(\theta \in [0,2\pi)\), with embedding
\[
	x(\theta) = (\cos\theta,\sin\theta) \in \mathbb R^2,
\]
and let \(\rho\) denote the uniform probability measure on \(S^1\),
\(d\rho(\theta)=d\theta/(2\pi)\). Given a limiting NTK
\(\bar\Theta(\theta,\theta')\), define the associated integral operator
\begin{equation}
	(T_{\bar\Theta}g)(\theta)
	:=
	\int_0^{2\pi}
	\bar\Theta(\theta,\theta')g(\theta')
	\,\frac{d\theta'}{2\pi}.
	\label{eq:bg-ntk-integral-operator}
\end{equation}
In the continuum constant-kernel regime, the residual
\(r_t(\theta)=f_t(\theta)-f^\star(\theta)\) evolves according to
\begin{equation}
	\partial_t r_t = -T_{\bar\Theta} r_t.
	\label{eq:bg-continuum-residual-dynamics}
\end{equation}
If \(\{\psi_j\}_{j\ge 0}\) is an orthonormal eigenbasis of
\(T_{\bar\Theta}\), with
\[
	T_{\bar\Theta}\psi_j=\lambda_j\psi_j,
\]
then the residual decomposes as
\[
	r_t = \sum_{j\ge 0} a_j(t)\psi_j,
	\qquad
	a_j(t)=\langle r_t,\psi_j\rangle_{L^2(\rho)}.
\]
Substituting this expansion into \eqref{eq:bg-continuum-residual-dynamics}
gives
\begin{equation}
	a_j(t)=e^{-\lambda_j t}a_j(0).
	\label{eq:bg-continuum-mode-decay}
\end{equation}
Thus each eigendirection is learned independently, and the decay rate is
determined by the corresponding eigenvalue. Larger eigenvalues correspond to
faster decay.

This is the operator form of the NTK training dynamics introduced by
\citet{jacot2018ntk}. In the context of spectral bias,
\citet{cao2020understandingspectralbiasdeep} use this eigenfunction viewpoint
to show that components aligned with larger NTK eigenvalues are learned faster.

For the circle, the Fourier basis appears when the limiting kernel is
rotation-invariant. This is the same symmetry principle used in analyses of NTK
spectra on the circle and sphere, where uniform input distributions lead to
Fourier modes on \(S^1\) and spherical harmonics on \(S^{d-1}\)
\citep{basri2019convergence,bietti2019inductivebiasneuraltangent,
	cao2020understandingspectralbiasdeep}. Under isotropic initialization, the distribution of \(w\)
is invariant under rotations. Using the infinite-width expression
\eqref{eq:ntk-1hl-infinite}, this implies
\[
	\bar\Theta(Rx,Rx')=\bar\Theta(x,x')
\]
for every rotation \(R\). Since
\(x(\theta+\alpha)=R_\alpha x(\theta)\), the kernel depends only on the angular
difference:
\begin{equation}
	\bar\Theta(\theta,\theta')
	=
	\kappa(\theta-\theta')
	\label{eq:bg-translation-invariant-kernel}
\end{equation}
for some \(2\pi\)-periodic function \(\kappa\).

With the uniform measure, \eqref{eq:bg-ntk-integral-operator} becomes a circular
convolution:
\begin{equation}
	(T_{\bar\Theta}g)(\theta)
	=
	\int_0^{2\pi}
	\kappa(\theta-\theta')g(\theta')
	\,\frac{d\theta'}{2\pi}
	=
	(\kappa*g)(\theta).
	\label{eq:bg-convolution-operator}
\end{equation}
Therefore the complex Fourier modes
\[
	\phi_q(\theta)=e^{iq\theta}, \qquad q\in\mathbb Z,
\]
are eigenfunctions of \(T_{\bar\Theta}\). Indeed,
\begin{align}
	(T_{\bar\Theta}\phi_q)(\theta)
	 & =
	\int_0^{2\pi}
	\kappa(u)e^{iq(\theta-u)}
	\,\frac{du}{2\pi} \\
	 & =
	e^{iq\theta}
	\int_0^{2\pi}
	\kappa(u)e^{-iqu}
	\,\frac{du}{2\pi}
	=
	\lambda_q \phi_q(\theta),
\end{align}
where
\begin{equation}
	\lambda_q
	=
	\int_0^{2\pi}
	\kappa(u)e^{-iqu}
	\,\frac{du}{2\pi}.
	\label{eq:bg-fourier-eigenvalues}
\end{equation}
When the kernel is real and even, one may equivalently use the real Fourier
basis
\[
	1,\qquad
	\sqrt{2}\cos(k\theta),\qquad
	\sqrt{2}\sin(k\theta),\qquad k\ge 1.
\]
In this basis, each Fourier frequency subspace evolves independently under the
continuum infinite-width operator.

This gives the basic spectral picture behind the NTK explanation of
frequency-dependent learning. In the idealized continuum infinite-width
setting, the Fourier modes are exact eigenfunctions of the training operator,
and their learning rates are determined by the corresponding eigenvalues
\(\lambda_q\). The next section briefly contrasts this fixed-operator view with
another large-width limit, before the thesis returns to the NTK framework.

\section{Another large-width viewpoint: mean-field}
\label{sec:background-mean-field-brief}

The previous section explains why the NTK viewpoint is useful for this thesis.
In the constant-kernel regime, training is described by a fixed operator on
functions. On \(S^1\), this operator is diagonal in the Fourier basis, so each
Fourier mode has a well-defined learning rate. This is the structure needed for
a mode-by-mode analysis of spectral bias.

There is another important large-width viewpoint, usually called the
\emph{mean-field} viewpoint. It starts from a different scaling of a two-layer
network,
\[
	f_\theta(x)
	=
	\frac{1}{m}
	\sum_{\alpha=1}^m
	v_\alpha\,\varphi(w_\alpha^\top x).
\]
Here each neuron has parameters
\[
	\theta_\alpha=(v_\alpha,w_\alpha)
	\in
	\Omega:=\mathbb R\times\mathbb R^d .
\]
Instead of following the \(m\) neurons one by one, the mean-field viewpoint
tracks how the neurons are distributed in parameter space. This is done through
the empirical measure
\[
	\mu_m
	:=
	\frac{1}{m}
	\sum_{\alpha=1}^m \delta_{\theta_\alpha},
\]
where \(\delta_{\theta_\alpha}\) places one unit of mass at the parameter value
\(\theta_\alpha\). With this notation, the network output can be written as
\[
	f_\theta(x)
	=
	\int_\Omega v\,\varphi(w^\top x)\,d\mu_m(v,w).
\]
Thus the network is represented by the distribution of its hidden units. During
training, this distribution moves through parameter space.

This gives a different limiting description from the NTK regime. In suitable
large-width limits, the empirical measure \(\mu_m\) converges to a limiting
time-dependent measure \(\mu_t\), and training is described by an evolution
equation for \(\mu_t\). This perspective was developed for two-layer networks
by \citet{Mei2018meanfieldview}, and related interacting-particle descriptions
were studied by \citet{Rotskoff_2022}. From an optimization viewpoint,
\citet{chizat2018globalconvergencegradientdescent} studied the corresponding
many-particle limit using tools from optimal transport. Extensions and rigorous
frameworks for deeper networks have also been studied
\citep{nguyen2019meanfieldlimitlearning,araujo2019meanfieldlimitcertaindeep,
	sirignano2019meanfieldanalysisneural,nguyen2023rigorousframeworkmeanfield}.

The reason we do not use this viewpoint as the main framework is not that it is
unrelated to neural-network training. It answers a different kind of question.
Mean-field limits are useful for describing how the hidden units move and how
the representation changes during training. In contrast, the spectral-bias
question studied here requires an operator acting on functions: we want to ask
which Fourier modes are eigenfunctions, what their eigenvalues are, and how the
corresponding eigenspaces behave. The NTK viewpoint gives such an operator
directly. The mean-field viewpoint describes the evolution of a parameter
distribution, so the same fixed Fourier-mode picture is not available in the
same direct form.

For this reason, the rest of the thesis uses the NTK formulation as the main
analytical framework. The mean-field derivation and its comparison with the NTK
scaling are given in Appendix~\ref{cha:mean-field-appendix}.
